\chapter{Structure of the Program}

\section{Program Entry Point}
Every EDADSL program must begin with the \texttt{start} keyword. The program executes sequentially from the entry point.

\section{Global Scope}
Variables and constants declared at the top level of the program operate on the global scope. They can be accessed from any point in the program after their declaration. However, variables declared inside functions are local to that function (see Section~\ref{sec:function-scope}).

\section{Execution Model}
% TODO: Define the formal execution model
% Uncertain about: 
% - How the program transitions between declaration phase and execution phase
% - Whether electrical statements are executed immediately or deferred
% - How the solver interacts with the program execution


\chapter{Structure of the Electrical System}

\section{Electrical Graph}
The electrical system is modeled as an undirected graph where:
\begin{itemize}
\item Vertices represent Nets
\item Edges represent Connections
\item Parts attach to the graph via their Pins
\end{itemize}

\section{Object Lifecycle}
% TODO: Define the lifecycle of electrical objects
% Uncertain about:
% - When objects are committed to the graph
% - How modifications to the graph are tracked
% - Undo/redo semantics if any


\chapter{Structure of the Solver}

\section{Constraint Representation}
Physical constraints are represented as optimization objectives with priorities.

\section{Resolution Algorithm}
% TODO: Define the formal resolution algorithm
% Uncertain about:
% - The specific optimization algorithm used
% - How constraints are weighted
% - Convergence criteria




\part{Hardware Model}
\chapter{Circuit Model}
The circuit model represents the electrical system as a graph of interconnected objects.

\section{Graph Structure}
The circuit is represented as a graph where:
\begin{itemize}
\item Nodes are Nets (electrical connection points)
\item Edges are Connections (physical wires/traces between pins)
\item Parts are components that connect to the graph via Pins
\end{itemize}

\section{Hierarchy}
The model supports hierarchical grouping:
\begin{itemize}
\item Parts contain Pins
\item Pins connect to Nets via Connections
\item Nets can be grouped by Tags
\end{itemize}


\chapter{Electrical Objects}
The electrical system consists of five primary object types.

\section{Nets}
A Net is a named electrical node that connects multiple Pins and Connections. Nets are the fundamental grouping mechanism for shared electrical connections.

\subsection{Creation}
Nets are created using the \texttt{elec.net} statement:
\begin{verbatim}
elec.net <name> [<tag1> ... <tag n>]

Example:
elec.net "Gnd"
elec.net "hvbus" "hv" "high_i" "noisy"
\end{verbatim}

\subsection{Properties}
\begin{itemize}
\item Name (unique identifier)
\item Connected objects (Pins and Connections)
\item Tags (user-defined labels)
\end{itemize}

\section{Parts}
A Part represents a physical electronic component in the circuit.

\subsection{Creation}
Parts are created using the \texttt{elec.part} statement:
\begin{verbatim}
elec.part <reference>: <model>(<args>); 
    des="<description>"; fp="<footprint>"

Example:
elec.part "R3": R(R=220, P=0.25, type="carbon"); 
    des="220 ohm resistor"; 
    fp="Resistor_THT:R_Axial_DIN0309"
\end{verbatim}

\subsection{Properties}
\begin{itemize}
\item Reference designator (e.g., R3, Q1, U1)
\item Model and parameters
\item Description
\item Footprint
\item Pins (derived from model/footprint)
\end{itemize}

\section{Pins}
A Pin is a connection point on a Part. Pins are accessed using the bracket notation.

\subsection{Creation}
Pins are created implicitly when Parts are instantiated.

\subsection{Reference Syntax}
\begin{verbatim}
<part_ref>[<pin_ref>]

Examples:
R1[1]    #c# Pin 1 of resistor R1
Q1[d]    #c# Drain pin of MOSFET Q1
U1[v+]   #c# Non-inverting input of op-amp U1
\end{verbatim}

\section{Connections}
A Connection is a binary electrical connection between exactly two Pins.

\subsection{Creation}
Connections can be created using the pipe syntax or \texttt{elec.connect}:
\begin{verbatim}
#c# Pipe syntax (creates a Connection object):
( R2[1] | J4[7] )

#c# Connect statement (joins pins to a Net):
elec.connect <net> <pin1> <pin2> ... <pin n>

Example:
elec.connect Gnd Q2[s] C6[-] R3[2]
\end{verbatim}

\section{Models}
Models define the electrical characteristics and available pins of a Part type.

\subsection{Model Declaration Syntax}
Custom models are defined using the following syntax:
\begin{verbatim}
EBNF:
model_declaration = "model", model_name, ":",
                    "{", pin_entries, [metadata], "}"

model_name        = id, "/", id

pin_entry         = pin_name, "->", pin_id, ";",
                    "T", "=", pin_type

pin_name          = id
pin_id            = integer
pin_type          = "input" | "power" | "output"

metadata          = des_entry | dts_entry | fp_entry

des_entry         = "des", "=", string_type
dts_entry         = "dts", "=", string_type
fp_entry          = "fp", "=", string_type
\end{verbatim}

The \texttt{model\_name} consists of a category and variant separated by \texttt{/} (e.g., \texttt{mosfet/C3M0016120K1}).

Pin types classify the electrical role:
\begin{itemize}
\item \texttt{input} -- Signal input (gate, base, control)
\item \texttt{power} -- Power connection (source, drain, emitter, collector)
\item \texttt{output} -- Signal output (drain for sense, kelvin connections)
\end{itemize}

\subsection{Built-in Models}
\begin{itemize}
\item \texttt{R} -- Resistor
\item \texttt{C} -- Capacitor
\item \texttt{L} -- Inductor
\item \texttt{opamp} -- Operational Amplifier
\item \texttt{MOSFET} -- MOSFET transistor
\item Additional models available in standard library
\end{itemize}

\subsection{Custom Model Examples}

\paragraph{MOSFET Model}
\begin{verbatim}
model MOSFET/IRLZ44N:{
    g -> 1; T = input
    S -> 2; T = power
    d -> 3; T = power
}
\end{verbatim}

\paragraph{Non-Polarized Capacitor}
\begin{verbatim}
model nonpol_C:{
    1 -> 1; T = power
    2 -> 2; T = power
}
\end{verbatim}

\paragraph{SiC MOSFET with Metadata}
\begin{verbatim}
model mosfet/C3M0016120K1:{
pins:
    g        -> 4; T = input
    s        -> 2; T = power
    d        -> 1; T = power
    kelvin_s -> 3; T = output

des: "1200 V, 125A, 16 mOhm, TO-247-4"

dts: "https://assets.wolfspeed.com/..."

fp:  "Package_TO_SOT_THT/TO-247-4_Vert"
}
\end{verbatim}

The \texttt{des} (description), \texttt{dts} (datasheet URL), and \texttt{fp} (footprint path) are optional metadata entries.

\section{Relationship between the Objects}
\begin{itemize}
\item Parts contain Pins
\item Pins are connected to Nets via Connections
\item Nets can have Tags for grouping
\item The global set \texttt{*e} contains all objects in creation order
\end{itemize}


\chapter{Querying}
EDADSL provides a powerful querying system for selecting electrical objects based on their relationships.

\section{Output and Input Quantifiers}
The quantifier syntax selects objects connected to a net:
\begin{verbatim}
<quantifier>@<net_name>
\end{verbatim}

\section{Net based Querying}
Quantifiers filter by object type:
\begin{verbatim}
P@<net>    #c# Parts connected to net
C@<net>    #c# Connections on net
T@<net>    #c# Pins connected to net
I@<net>    #c# Connections and Pins on net
H@<net>    #c# Holes on net
E@<net>    #c# Everything on net (or just @<net>)
\end{verbatim}

Example:
\begin{verbatim}
$input [set]= I@Vin
\end{verbatim}

\section{Tagged Net based Querying}
Query by tag instead of net name:
\begin{verbatim}
<quantifier>@^<tag>
\end{verbatim}
This expands to the union of the quantifier across all nets with the given tag.

Example:
\begin{verbatim}
elec.tag Vin "sensitive"
elec.tag Vfb "sensitive"
I@^sensitive  #c# All connections/pins on 
              #c# nets tagged "sensitive"
\end{verbatim}

\section{General Querying}
\label{sec:general-querying}
General querying provides two powerful forms for filtering and transforming sets: set builder notation and set query notation. Both support nesting.

\subsection{Set Builder Notation}
Set builder notation iterates over a set or list, filtering and optionally transforming elements.
\begin{verbatim}
EBNF:
set_builder = "{", variable, "E", ( set | list ), 
              ":", filter_expr, [ "~", transform_expr ], "}"

filter_expr   = function_call   #c# returns bool
transform_expr = function_call  #c# applied to filtered elements
\end{verbatim}

The filter function receives each element and returns \texttt{True} to keep it. The transform function is applied \textbf{per-element} to each item that passes the filter; the result set contains the transformed values. If transform is omitted, elements are returned as-is.

Examples:
\begin{verbatim}
#c# Filter: only active parts
{ $x E *p : f.u.is_active($x) ~ noop($x) }

#c# Filter and transform: get names of large caps
{ $x E *p : f.u.is_capacitor($x) && 
             f.u.value($x) > 100u ~ 
             f.u.name($x) }

#c# Nested: find all nets connected to parts 
#c# in a set
{ $x E *p : f.u.is_mosfet($x) ~ 
    { $y E *n : f.u.connected($y, $x) ~ 
        noop($y) } }
\end{verbatim}

\subsection{Set Query Notation}
Set query notation applies a filter and optional transform to an existing set, with optional output assignment.
\begin{verbatim}
EBNF:
set_query = "{", set_expr, "@~", "[", 
            filter_expr, "]", 
            [ "~[", transform_expr, "]" ],
            "~>", [ variable ], "}"
\end{verbatim}

The \texttt{@\~} operator queries the set. The trailing \texttt{\~{}>} optionally writes the result to a variable; if omitted, the result is returned directly.

Examples:
\begin{verbatim}
#c# Query with filter only
{ *p @~[f.u.is_active] ~> }

#c# Query with filter and transform
{ *p @~[f.u.is_mosfet] 
      ~[f.u.name] ~> }

#c# Store result in variable
{ *p @~[f.u.is_capacitor] 
      ~[f.u.value] ~> $cap_values }

#c# Query connections on a net
{ C@Vin @~[f.u.is_high_current] ~> }
\end{verbatim}

\subsection{Nesting}
Both forms can be nested to any depth. Each nested query creates a new scope for its loop variable.

Example:
\begin{verbatim}
#c# For each MOSFET, find all connected nets,
#c# then filter to high-voltage nets
{ $q E *p : f.u.is_mosfet($q) ~ 
    { C@^power @~[f.u.connected_to($q)] 
               ~[f.u.voltage_rating] ~> 
      $ratings } }
\end{verbatim}\part{Physical System}
